\section{Safety, Oversight, and Ethical Considerations}
\label{sec:safety}

The DGM authors are explicit that their experiments relied on sandboxing and
human oversight, and they note that the broader safety challenges of recursively
self-improving systems remain open~\cite{zhang2025dgm}. SEAL's authors
similarly emphasize that human supervision is required to ensure improvements
remain beneficial and aligned~\cite{zweiger2025seal}. Evolving educational
content raises its own distinct concerns that demand comparable care.

\textbf{Goodhart's Law and proxy gaming.} Any fitness function is a proxy for
genuine learning, and optimizing a proxy can degrade the true objective. A
variant that maximizes quiz scores by narrowing instruction to test items, or
that maximizes engagement through manipulative design, would score well while
harming learning. The multi-dimensional fitness function of
Section~\ref{sec:fitness}, with heavy weight on retention and normalized gain
rather than raw scores, is a partial defense, but human review of \textit{why}
a variant performs well remains essential.

\textbf{Human-in-the-loop gating.} No generated variant should reach learners
without review. The \texttt{human\_review} block in the fitness schema enforces
this structurally: a variant cannot become \texttt{active} without an approving
reviewer. This mirrors the DGM's fixed, human-controlled parent-selection and
archiving process, which the authors describe as an implicit safety constraint.

\textbf{Equity across learner segments.} Because the framework supports
audience-conditioned evolution, it risks optimizing well for majority segments
while neglecting smaller ones. Fitness aggregation should monitor per-segment
performance and flag variants that improve aggregate scores while harming any
segment, particularly underserved populations such as rural or under-resourced
learners.

\textbf{Privacy.} The schema deliberately uses account-based identifiers rather
than personal information and supports a fully anonymous mode. Learner data
feeding the evolutionary loop should be governed by clear consent and
data-minimization practices, and identifiers should never be embedded in URLs
or query strings.

\textbf{Sandboxed deployment.} Generated MicroSim code should be validated in
an isolated preview environment before any public deployment, a practice that
aligns naturally with a preview-branch workflow on a static host.
